\documentclass[sigconf,nonacm]{acmart}

\usepackage[T1]{fontenc}
\usepackage{booktabs}
\usepackage{subcaption}
\usepackage{enumitem}
\usepackage{algorithm}
\usepackage{algpseudocode}

\settopmatter{printacmref=false}
\renewcommand\footnotetextcopyrightpermission[1]{}
\pagestyle{plain}

\acmConference[ETRA '27]{2027 Symposium on Eye Tracking Research and Applications}{May 2027}{Stuttgart, Germany}
\acmYear{2027}
\copyrightyear{2027}

\title{OpenGaze: Head-Invariant In-Browser Gaze Estimation with In-Situ Online Fine-Tuning and Zero-DOM Bayesian Reading Telemetry}

\author{Rahul Drabit}
\affiliation{%
  \institution{Cognitive Vision and Interaction Lab}
  \country{Independent Researcher}
}
\email{contact@rahuldrabit.dev}

\begin{abstract}
Commodity webcam-based eye tracking holds immense potential to democratize behavioral psychology, reading assessment, cognitive diagnostics, and accessible human-computer interaction by eliminating the need for expensive, proprietary infrared eye trackers. However, broad ecological adoption has been severely impeded by four compounding failure modes: (1) acute vulnerability to natural user head sway, where as little as a single centimeter of involuntary lateral head translation produces feature displacements in camera pixel space that exceed the biological eye rotation across an entire 24-inch display; (2) numerical ill-conditioning in classical polynomial calibration mappings ($\kappa(A^T A) > 10^5$), which triggers catastrophic overfitting, runaway regression coefficients, and severe gaze jitter; (3) rapid temporal posture decay (drift), where calibration degrades within three to five minutes as participants settle into their chairs; and (4) intrusive Document Object Model (DOM) mutation—such as wrapping paragraph text into thousands of individual \texttt{<span>} elements—which degrades browser rendering performance, alters typographic line flow, and invalidates natural reading dynamics.

In this paper, we present \textbf{OpenGaze}, an open-source, mathematically principled in-browser eye-tracking framework engineered to resolve these challenges in commodity web environments without specialized hardware or invasive page alteration. OpenGaze introduces four primary contributions: (1) a closed-form \textit{canthal vector projection} that anchors iris displacements to rigid anatomical skull landmarks, decoupling 1st-order head translation, distance scaling, and roll tilt in camera pixel space; (2) a 28-term physics-informed feature space incorporating perspective foreshortening, metric depth proxies ($Z_{\text{ratio}}$), and eyelid occlusion compensation, solved via multi-task \textit{ElasticNet coordinate descent}; (3) an \textit{in-situ continuous fine-tuning engine} that opportunistically harvests passive eye-hand foveation pairs during natural click interactions and anchors updates to the baseline calibration via an elastic Bayesian prior to prevent catastrophic spatial divergence; and (4) a \textit{zero-DOM hierarchical Bayesian reading telemetry engine} that queries native browser layout rectangles via \texttt{Range.getClientRects()} and incorporates human Preferred Viewing Location (PVL = 0.40) priors to achieve robust sub-line word attribution.

We empirically validate OpenGaze against four competitive baseline architectures across multi-session in-browser datasets ($N=3$ independent recording sessions, 1,045 captured video frames, 20-target serpentine grid) under strict held-out 20-fold Grouped Leave-One-Target-Out Cross-Validation. OpenGaze achieves a Mean Cross-Validation Root Mean Square (CV RMS) error of \textbf{138.8\,px} across all sessions (reaching \textbf{101.0\,px} on optimal sessions, an 80.6\% error reduction over unregularized WebGazer-style polynomial baselines) with an average inference latency of only \textbf{0.837\,ms} per frame on a single commodity CPU thread. Systematic ablation studies demonstrate that discarding saccadic transit latency frames reduces calibration error inflation by 47.1\%, while incorporating the human PVL landing prior boosts exact word attribution accuracy from 68.1\% to \textbf{89.4\%}. The complete framework, diagnostic test suite, and benchmark runner are released under the MIT license at \url{https://github.com/Rahuldrabit/Eye_Tracker}.
\end{abstract}

\begin{CCSXML}
<ccs2012>
   <concept>
       <concept_id>10003120.10003121.10003128</concept_id>
       <concept_desc>Human-centered computing~Interaction techniques</concept_desc>
       <concept_significance>500</concept_significance>
   </concept>
   <concept>
       <concept_id>10010147.10010178.10010224</concept_id>
       <concept_desc>Computing methodologies~Computer vision</concept_desc>
       <concept_significance>500</concept_significance>
   </concept>
   <concept>
       <concept_id>10003120.10003121.10003122</concept_id>
       <concept_desc>Human-centered computing~HCI design and evaluation methods</concept_desc>
       <concept_significance>300</concept_significance>
   </concept>
</ccs2012>
\end{CCSXML}

\ccsdesc[500]{Human-centered computing~Interaction techniques}
\ccsdesc[500]{Computing methodologies~Computer vision}
\ccsdesc[300]{Human-centered computing~HCI design and evaluation methods}

\keywords{Webcam eye tracking, gaze estimation, head-pose invariance, online adaptation, Bayesian attribution, reading telemetry, ElasticNet, Preferred Viewing Location, in-browser computer vision.}

\maketitle

\section{Introduction}
\label{sec:intro}
Eye tracking represents one of the most powerful behavioral windows into human cognitive processing, attentional allocation, reading comprehension, and neurological function \cite{rayner1998eye,holmqvist2011eye,reichle2003ez}. For decades, experimental eye tracking has been dominated by specialized Pupil Center Corneal Reflection (PCCR) hardware (e.g., EyeLink, Tobii, SMI). While these commercial instruments achieve exceptional spatial accuracy ($< 0.5^\circ$) and high sampling rates (250--2000\,Hz), their prohibitive cost (\$10,000--\$45,000), requirement for dedicated hardware setups, and physical laboratory confinement severely restrict their utility for ecological, large-scale studies \cite{papoutsaki2016webgazer,krafka2016eye}.

The advent of client-side web technologies and real-time computer vision libraries—most notably WebGazer.js \cite{papoutsaki2016webgazer} and Google MediaPipe FaceMesh \cite{lugaresi2019mediapipe}—sparked widespread enthusiasm for webcam-based gaze tracking. By executing directly within commodity web browsers on standard consumer laptops, these systems hold the promise of democratizing visual attention metrics for education, digital reading assessment, e-health, and accessible computing across millions of users without specialized hardware.

Despite more than a decade of research, however, commodity webcam gaze tracking has struggled to transition from coarse quadrant selection into fine-grained behavioral measurement (such as determining which word or line of text a user is reading). In unconstrained real-world environments, consumer webcam gaze estimation is undermined by four foundational, interlocking failure modes:

\begin{enumerate}[leftmargin=*]
    \item \textbf{Head-Movement Signal Dominance}: Monocular webcams lack the active infrared illumination required to generate corneal reflections (glints). Consequently, software trackers must infer gaze from iris landmark displacements in the camera's 2D image plane. When landmark coordinates are normalized to the video frame bounding box ($[0, 1] \times [0, 1]$), biological eye rotation across a typical 24-inch monitor at a 60\,cm viewing distance displaces the iris center by approximately 0.0125 normalized units. In contrast, natural involuntary head sway (respiration, postural settling, or typing vibration) of just 10\,mm translates the same landmark by approximately 0.0175 normalized units \cite{luger2020invariance}. Under standard webcam architectures, natural head movement does not merely perturb the gaze signal—it completely overwhelms it, causing catastrophic calibration failure.
    
    \item \textbf{Ill-Conditioned Calibration Mappings}: Classical in-browser frameworks map eye coordinates to screen pixel targets by fitting second-order polynomials ($[1, x, y, xy, x^2, y^2]$) via unregularized Ordinary Least Squares (OLS) \cite{papoutsaki2016webgazer,blignaut2009mapping}. Because normalized iris coordinates occupy an extremely narrow numeric domain ($x, y \approx 0.50 \pm 0.04$), quadratic powers ($x^2, xy$) are near-linear combinations of the linear terms. The resulting Gram matrix $A^T A$ exhibits severe numerical rank-deficiency, with condition numbers routinely exceeding $\kappa(A^T A) > 10^5$. This ill-conditioning inflates regression weights to unstable magnitudes ($10^4$--$10^6$), transforming sub-pixel landmark jitter into tens of pixels of erratic screen scatter.
    
    \item \textbf{Temporal Posture Decay (Drift)}: Standard eye-tracking protocols treat calibration as a static, one-time ceremony conducted at the beginning of a recording session. However, empirical studies demonstrate that calibration quality deteriorates rapidly within 3 to 5 minutes as users naturally slump, reposition their seating, or experience ambient lighting variations \cite{blignaut2014mapping}. Existing web trackers lack self-stabilizing mechanisms to absorb this postural drift online without interrupting the user with intrusive recalibration dialogs.
    
    \item \textbf{Intrusive DOM Mutation in Reading Telemetry}: In educational and psychological reading research, mapping gaze fixations to specific textual elements has traditionally required mutating the host page by wrapping every word into an individual Document Object Model (DOM) node (e.g., \texttt{<span data-word="i">}). On an average reading passage of 1,000 words, this practice multiplies the document node count tenfold, damages CSS layout fidelity, disrupts native text selection, and forces continuous, expensive layout reflows during 60\,FPS gaze rendering.
\end{enumerate}

To overcome these fundamental barriers, this paper presents \textbf{OpenGaze}, an open-source, modular, and mathematically principled web eye-tracking framework. OpenGaze bridges the gap between commodity consumer webcams and research-grade reading telemetry through five core innovations:

\begin{itemize}[leftmargin=*]
    \item \textbf{Closed-Form Canthal Vector Projection}: An intrinsic 2D coordinate normalization anchored to rigid skull canthi ($p_{\text{inner}}, p_{\text{outer}}$) that projects iris offsets into an orthonormal, skull-relative basis in camera pixel space, proving exact 1st-order invariance to head translation, distance scaling, and roll tilt.
    \item \textbf{28-Term Physics-Informed ElasticNet Solver}: A structured feature space capturing perspective foreshortening, optical depth proxies ($Z_{\text{ratio}}$), binocular disparity, and eyelid occlusion, optimized via standardized ElasticNet coordinate descent that enforces $\ell_1$ sparsity on spurious cross-terms while maintaining $\ell_2$ grouping.
    \item \textbf{In-Situ Continuous Fine-Tuning Engine}: A self-supervised online adaptation mechanism that extracts passive eye-hand foveation pairs during natural mouse clicks ($t \in [t_c - 240\,\text{ms}, t_c - 60\,\text{ms}]$) and anchors weight updates to initial calibration via an elastic Bayesian prior to prevent spatial distortion.
    \item \textbf{Zero-DOM Hierarchical Bayesian Reading Telemetry}: A completely non-invasive reading telemetry engine that extracts word bounding geometries directly from native browser text nodes via \texttt{Range.getClientRects()} and applies a two-stage Bayesian posterior incorporating the human Preferred Viewing Location (PVL = 0.40) prior.
    \item \textbf{Extensive Empirical Validation}: A comprehensive benchmark across multi-session in-browser datasets ($N=3$ sessions, 1,045 frames) evaluating five baseline architectures under strict 20-fold Grouped Leave-One-Target-Out Cross-Validation, demonstrating an 80.6\% error reduction over classical WebGazer baselines at an inference latency of $0.837\,$ms.
\end{itemize}

The remainder of this paper is organized as follows: Section~\ref{sec:related} surveys related literature. Section~\ref{sec:math} details the mathematical architecture of OpenGaze. Section~\ref{sec:online} formalizes the in-situ online adaptation engine. Section~\ref{sec:bayesian} introduces the zero-DOM Bayesian text telemetry model. Section~\ref{sec:experiments} presents the empirical evaluation against baselines. Section~\ref{sec:ablation} provides systematic ablation experiments. Section~\ref{sec:discussion} discusses practical deployment guidelines, limitations, and ethical safeguards. Finally, Section~\ref{sec:conclusion} concludes and outlines future work.

\section{Related Work}
\label{sec:related}

\subsection{Appearance-Based vs. Geometric Gaze Estimation}
Gaze estimation methods are broadly categorized into appearance-based deep learning and geometric feature regression \cite{hansen2010eye,holmqvist2011eye}. Appearance-based methods feed raw eye or face image patches into convolutional neural networks (CNNs) or vision transformers trained on extensive crowdsourced datasets, such as MPIIGaze \cite{zhang2015appearance}, GazeCapture \cite{krafka2016eye}, and TurkerGaze \cite{xu2015turkergaze}. While these models demonstrate robust generalizability across diverse ethnicities and lighting conditions, their practical deployment within client-side web browsers remains severely constrained. Multi-layer neural backbones executing via WebGL or WebAssembly (WASM) consume substantial GPU memory bandwidth and induce thermal throttling on battery-powered mobile devices and low-cost laptops, leaving inadequate computational headroom for host educational or experimental applications.

Consequently, lightweight landmark-based geometric regression remains the gold standard for responsive in-browser interaction. WebGazer.js \cite{papoutsaki2016webgazer} pioneered this space by demonstrating that pupil center landmarks tracked in real-time could be mapped to screen coordinates via local regression models trained on mouse clicks. Similarly, Google MediaPipe FaceMesh \cite{lugaresi2019mediapipe} introduced a high-performance 468-point 3D facial mesh augmented with 10 dedicated iris rim landmarks executing at $>60\,$FPS in WebAssembly. However, raw landmark trackers output coordinates relative to the video frame viewport. Without explicit skull-relative normalization, user head movements corrupt landmark positions, causing severe gaze misestimation \cite{luger2020invariance}.

\subsection{Head-Pose Invariance and Normalization}
The vulnerability of gaze trackers to head motion has been widely documented in ocular literature \cite{blignaut2009mapping,blignaut2014mapping,sesin2008investigation}. In commercial infrared systems, head pose is canceled geometrically because infrared corneal glints move synchronously with the skull, allowing the vector difference between pupil center and glint to cancel translation \cite{hansen2010eye}. In passive webcam tracking, however, no glint signal exists.

Several computer vision techniques attempt head pose compensation by estimating full 3D head orientation ($\text{yaw}, \text{pitch}, \text{roll}$) via Perspective-n-Point (PnP) solvers aligned to an idealized 3D anthropometric face model \cite{sugano2014learning,zhang2015appearance}. While theoretically sound, 3D head pose estimation from monocular consumer webcams is notoriously sensitive to lens distortion, camera sensor noise, and individual facial morphology variations. A residual error of just $1.5^\circ$ in estimated pitch or yaw maps to an error of over 45 pixels on a 1080p display at a 60\,cm viewing distance. In contrast, OpenGaze circumvents non-linear 3D mesh optimization entirely, formulating an exact, closed-form 2D projection anchored to rigid anatomical canthal landmarks.

\subsection{Temporal Filtering and Fixation Detection}
Raw gaze coordinate streams from webcams exhibit two distinct noise profiles: high-frequency tremor (induced by landmark jitter and camera quantization noise) and low-frequency drift (induced by user posture decay) \cite{holmqvist2011eye}. Standard low-pass filters, such as Exponential Moving Average (EMA) or Gaussian smoothing, face an intractable trade-off: sufficient smoothing to eliminate fixation jitter introduces severe phase lag during rapid ballistic saccades (which reach velocities of 300--800$^\circ$/s), smearing gaze points across unrelated screen regions \cite{casiez20121}.

Casiez et al. \cite{casiez20121} resolved this latency-jitter trade-off by introducing the 1-Euro filter, an adaptive low-pass filter whose cutoff frequency dynamically scales with derivative signal velocity. In fixation identification, Salvucci and Goldberg's seminal taxonomy \cite{salvucci2000identifying} contrasted Velocity-Threshold Identification (I-VT) and Dispersion-Threshold Identification (I-DT). While I-VT is optimal for high-speed eye trackers ($\ge 250\,$Hz), I-DT is markedly more robust for lower-frequency webcam streams ($30\text{--}60\,$Hz), where numerical differentiation of landmark coordinates amplifies sensor noise. OpenGaze builds upon these foundations by coupling an order-gated 1-Euro filter with an anisotropic, uncertainty-scaled I-DT detector that enforces a seed-forward saccade transition architecture.

\subsection{Reading Telemetry and Preferred Viewing Location}
During silent reading, human eye movements do not sweep smoothly across text; instead, they alternate between stationary fixations (typically lasting 200--250\,ms) and rapid ballistic saccades spanning 7--9 character spaces \cite{rayner1998eye,reichle2003ez}. Landmark psycho-linguistic investigations by Rayner \cite{rayner1979preferred} and McConkie and Rayner \cite{mcconkie1979initial} established the existence of the human \textit{Preferred Viewing Location} (PVL): in left-to-right alphabetic orthographies, initial foveal fixations do not land at the geometric midpoint of a word, but systematically cluster between 0.35 and 0.40 of the word's width (measured from the left edge).

Existing digital reading research tools frequently neglect this physiological landing bias, naively mapping fixation centroids to the nearest word bounding box center ($0.50 \cdot \text{width}$). When coupled with webcam calibration noise, this midpoint assumption produces a systematic rightward attribution bias, misattributing fixations to subsequent words. OpenGaze directly incorporates the PVL prior into a hierarchical Bayesian formulation while querying raw layout rectangles via native browser APIs, completely eliminating the need for invasive DOM span-wrapping.

\section{Mathematical Architecture}
\label{sec:math}

The complete OpenGaze processing pipeline operates as a deterministic, decoupled feed-forward system, as illustrated in Figure~\ref{fig:pipeline}: Video Frame Capture $\rightarrow$ Canthal Basis Projection $\rightarrow$ 28-Term Feature Expansion $\rightarrow$ ElasticNet Screen Mapping $\rightarrow$ Order-Gated 1-Euro Filtering $\rightarrow$ Anisotropic I-DT Fixation Extraction $\rightarrow$ Zero-DOM Bayesian Text Attribution.

\begin{figure*}[t]
\centering
\fbox{\parbox{0.96\textwidth}{
\centering
\textbf{OpenGaze Core Mathematical Processing Pipeline} \\[4pt]
\small
$\text{Webcam Frame (720p)} \xrightarrow{\text{MediaPipe}} \text{Landmarks } (p_{\text{iris}}, p_{\text{inner}}, p_{\text{outer}}) \xrightarrow{\text{Section } \ref{subsec:canthal}} \text{Canthal Basis } (g_x, g_y) \xrightarrow{\text{Section } \ref{subsec:features28}} \text{28-D Vector } \boldsymbol{\phi}$ \\[3pt]
$\boldsymbol{\phi} \xrightarrow{\text{Section } \ref{subsec:elasticnet}} \text{ElasticNet Solver} \xrightarrow{\text{Section } \ref{subsec:filter}} \text{1-Euro Filter} \xrightarrow{\text{Section } \ref{subsec:fixation}} \text{Uncertainty I-DT} \xrightarrow{\text{Section } \ref{sec:bayesian}} \text{Bayesian Word Attribution (PVL)}$
}}
\caption{Architectural overview of the OpenGaze pipeline showing the end-to-end transformation from raw webcam frames to Bayesian word-level reading telemetry.}
\label{fig:pipeline}
\end{figure*}

\subsection{Head-Invariant Canthal Vector Projection}
\label{subsec:canthal}
To achieve invariance to 1st-order head translation, distance scaling, and roll tilt without requiring iterative 3D head pose estimation, OpenGaze anchors eye movements to rigid anatomical facial landmarks: the medial (inner) canthus ($p_{\text{inner}} \in \mathbb{R}^2$) and the lateral (outer) canthus ($p_{\text{outer}} \in \mathbb{R}^2$). These canthal points are anatomically anchored to the lacrimal bone and zygomatic frontal process of the human skull, exhibiting negligible displacement during facial expressions or eye movements.

Crucially, all landmark coordinates must first be projected from the camera's normalized texture space ($[0, 1] \times [0, 1]$) into physical camera pixel space:
\begin{equation}
    p_x = \tilde{p}_x \cdot W_{\text{frame}}, \quad p_y = \tilde{p}_y \cdot H_{\text{frame}}
\end{equation}
Because webcams capture frames at non-unity aspect ratios (typically 16:9 or 4:3, where $W_{\text{frame}} \neq H_{\text{frame}}$), operating directly on normalized texture coordinates distorts Euclidean distances, rendering the horizontal and vertical eye axes non-orthogonal.

For each individual eye, we define the intrinsic anatomical coordinate system:
\begin{align}
    m &= \frac{p_{\text{inner}} + p_{\text{outer}}}{2} \quad \text{(Canthal Midpoint)} \\
    w &= \|p_{\text{outer}} - p_{\text{inner}}\|_2 \quad \text{(Inter-canthal Width)} \\
    \mathbf{u} &= s \cdot \frac{p_{\text{outer}} - p_{\text{inner}}}{w}, \quad s = {\begin{cases} -1 & \text{left eye (landmark 468)} \\ +1 & \text{right eye (landmark 473)} \end{cases}} \\
    \mathbf{v} &= \begin{pmatrix} -u_y \\ u_x \end{pmatrix} \quad \text{(Orthonormal Eye Vertical Axis)}
\end{align}

The directional sign multiplier $s$ accounts for the bilateral asymmetry of human facial anatomy: for the left eye (image-left), the lateral outer corner lies to the left of the medial inner corner, whereas for the right eye (image-right), the outer corner lies to the right. Without the sign inversion $s$, the basis vectors of the two eyes point in opposite directions, causing binocular averaging to cancel the gaze signal rather than reinforce it.

Let $p_{\text{iris}} \in \mathbb{R}^2$ denote the iris center, calculated as the centroid of the center landmark and the four boundary rim landmarks ($468 + 469\text{--}472$ for left; $473 + 474\text{--}477$ for right). The projected, dimensionless gaze offset $\mathbf{g} = (g_x, g_y)^T$ in units of eye-widths is:
\begin{equation}
    g_x = \frac{(p_{\text{iris}} - m) \cdot \mathbf{u}}{w}, \quad g_y = \frac{(p_{\text{iris}} - m) \cdot \mathbf{v}}{w}
\end{equation}

\begin{theorem}[Head Invariance of Canthal Basis]
Under uniform rigid skull translation $\Delta \mathbf{t} \in \mathbb{R}^2$, scale magnification $\kappa > 0$, and in-plane roll rotation $\mathbf{R}(\theta) \in \text{SO}(2)$, the normalized canthal gaze coordinates $\mathbf{g}$ remain perfectly invariant to first order.
\end{theorem}

\begin{proof}
Consider a rigid planar transformation of the skull given by $p' = \kappa \mathbf{R}(\theta) p + \Delta \mathbf{t}$.
First, under translation $\Delta \mathbf{t}$:
\begin{equation}
    p_{\text{iris}}' - m' = (p_{\text{iris}} + \Delta \mathbf{t}) - (m + \Delta \mathbf{t}) = p_{\text{iris}} - m
\end{equation}
Second, the inter-canthal width transforms as:
\begin{equation}
    w' = \|p_{\text{outer}}' - p_{\text{inner}}'\|_2 = \kappa \|\mathbf{R}(\theta)(p_{\text{outer}} - p_{\text{inner}})\|_2 = \kappa w
\end{equation}
Third, the basis vector $\mathbf{u}'$ transforms as:
\begin{equation}
    \mathbf{u}' = s \cdot \frac{\kappa \mathbf{R}(\theta)(p_{\text{outer}} - p_{\text{inner}})}{\kappa w} = \mathbf{R}(\theta) \mathbf{u}
\end{equation}
The orthonormal vertical vector transforms identically: $\mathbf{v}' = \mathbf{R}(\theta) \mathbf{v}$.
Finally, computing the projected coordinate $g_x'$:
\begin{equation}
    g_x' = \frac{(p_{\text{iris}}' - m') \cdot \mathbf{u}'}{w'} = \frac{\kappa \mathbf{R}(\theta)(p_{\text{iris}} - m) \cdot \mathbf{R}(\theta)\mathbf{u}}{\kappa w}
\end{equation}
Because the rotation matrix $\mathbf{R}(\theta)$ is unitary orthogonal, the inner product is preserved: $(\mathbf{R} \mathbf{a}) \cdot (\mathbf{R} \mathbf{b}) = \mathbf{a} \cdot \mathbf{b}$. The scaling factor $\kappa$ cancels exactly in the numerator and denominator:
\begin{equation}
    g_x' = \frac{\mathbf{a} \cdot \mathbf{b}}{w} = g_x
\end{equation}
An identical identity holds for $g_y'$. Thus, the dimensionless canthal gaze coordinates $\mathbf{g} = (g_x, g_y)^T$ are completely invariant to 2D translation, distance scaling, and roll tilt.
\end{proof}

\subsection{28-Term Physics-Informed Feature Space}
\label{subsec:features28}
While the canthal projection eliminates 1st-order planar head movement, 2nd-order non-linear physical phenomena remain: perspective foreshortening as the face rotates in 3D, camera distance fluctuations, binocular vergence, and eyelid occlusion during downward gaze. OpenGaze captures these phenomena by expanding the base canthal offsets into a structured 28-dimensional feature vector $\boldsymbol{\phi} \in \mathbb{R}^{28}$:

\begin{enumerate}[leftmargin=*]
    \item \textbf{Intercept and Base Dimensionless Gaze ($\phi_{0:2}$)}: $\phi_0 = 1.0$, $\phi_1 = g_x$, $\phi_2 = g_y$. Here $g_x, g_y$ represent the binocular average of the left and right eyes.
    \item \textbf{Independent Eye Channels ($\phi_{3:6}$)}: Asymmetric offsets $\phi_{3:6} = [g_{x,L}, g_{y,L}, g_{x,R}, g_{y,R}]$. This models ocular dominance, strabismus, and asymmetric illumination effects.
    \item \textbf{Head-Pose Anchors ($\phi_{7:10}$)}: Dimensionless yaw and pitch proxies derived from rigid facial landmarks:
    \begin{equation}
        \text{yawProxy} = \frac{d_{\text{cheek},L} - d_{\text{cheek},R}}{d_{\text{cheek},L} + d_{\text{cheek},R}}, \quad \text{pitchProxy} = \frac{p_{\text{nose},y} - p_{\text{forehead},y}}{\text{faceScale}}
    \end{equation}
    where $d_{\text{cheek},L}, d_{\text{cheek},R}$ are Euclidean distances from the nose bridge to the outer jawline, and $\text{faceScale}$ is the inter-pupillary baseline distance.
    \item \textbf{Metric Optical Depth Ratio ($\phi_{11:12}$)}:
    \begin{equation}
        Z_{\text{ratio}} = \frac{1}{2} \left(\frac{\text{faceScale}}{\text{faceScale}_{\text{anchor}}} + \frac{r_{\text{iris}}}{r_{\text{anchor}}}\right) - 1.0
    \end{equation}
    where $\text{faceScale}_{\text{anchor}}$ and $r_{\text{anchor}}$ are the geometric baselines established during calibration target 0. When the user leans closer to the screen, $Z_{\text{ratio}} > 0$; when leaning backward, $Z_{\text{ratio}} < 0$.
    \item \textbf{Parallax and Depth Coupling ($\phi_{13:16}$)}: First-order interaction cross-products between gaze orientation and depth: $\phi_{13:16} = [g_x \cdot Z_{\text{ratio}}, g_y \cdot Z_{\text{ratio}}, \text{pitch} \cdot Z_{\text{ratio}}, \text{yaw} \cdot Z_{\text{ratio}}]$.
    \item \textbf{2nd-Order Geometric Polynomials ($\phi_{17:20}$)}: Quadratic terms $\phi_{17:20} = [g_x^2, g_y^2, g_x \cdot g_y, g_x^2 + g_y^2]$ capturing spherical eye rotation mapped to a planar monitor surface.
    \item \textbf{Eyelid Occlusion Compensation ($\phi_{21:24}$)}: During reading, reading lower text lines requires eye depression, causing the upper eyelid to partially occlude the superior iris rim. We model this via the Eye Aspect Ratio ($\text{EAR} = \frac{\|p_{\text{top}} - p_{\text{bottom}}\|}{2 \|p_{\text{inner}} - p_{\text{outer}}\|} $), setting $\phi_{21:24} = [g_y \cdot \text{EAR}_L, g_y \cdot \text{EAR}_R, \text{EAR}_L, \text{EAR}_R]$.
    \item \textbf{Screen Tangent Projections ($\phi_{25:27}$)}: Cubic non-linear terms $\phi_{25:27} = [g_x^3, g_y^3, g_x \cdot g_y^2]$, representing the planar tangent distortion $\tan(\theta) \approx \theta + \frac{1}{3}\theta^3$ that arises when projecting angular gaze sweeps onto flat computer displays.
\end{enumerate}

\subsection{Standardized Multi-Task ElasticNet Solver}
\label{subsec:elasticnet}
Directly fitting screen coordinates $\mathbf{y} \in \mathbb{R}^n$ from a 28-dimensional collinear feature space using Ordinary Least Squares leads to severe overfitting. To ensure numerical stability, all feature columns ($j = 1 \dots 27$) are standardized across calibration samples to zero mean and unit variance:
\begin{equation}
    \mu_j = \frac{1}{n} \sum_{i=1}^n \phi_{i,j}, \quad \sigma_j = \sqrt{\frac{1}{n} \sum_{i=1}^n (\phi_{i,j} - \mu_j)^2 + \epsilon}, \quad Z_{i,j} = \frac{\phi_{i,j} - \mu_j}{\sigma_j}
\end{equation}
where $\epsilon = 10^{-8}$ prevents division by zero for invariant terms. The intercept column $Z_{i,0} = 1.0$ is left unstandardized and is never penalized.

We formulate the calibration mapping as an ElasticNet optimization problem \cite{zou2005regularization}, solved independently for the horizontal screen axis $y_x$ and vertical axis $y_y$:
\begin{equation}
    \min_{\mathbf{w} \in \mathbb{R}^{28}} \frac{1}{2n} \|\mathbf{y} - Z \mathbf{w}\|_2^2 + \alpha \rho \|\mathbf{w}_{1:27}\|_1 + \frac{1}{2} \alpha (1 - \rho) \|\mathbf{w}_{1:27}\|_2^2
\end{equation}
where $\alpha > 0$ controls overall regularization shrinkage, and $\rho \in [0, 1]$ governs the ElasticNet mixing ratio. We set $\rho = 0.30$: this strongly leverages $\ell_2$ ridge shrinkage to group collinear depth and parallax features, while maintaining sufficient $\ell_1$ sparsity to drive redundant cubic cross-terms to zero.

The objective is solved using cyclical coordinate descent. For each feature coordinate $j = 0, \dots, 27$, we compute the partial residual correlation:
\begin{equation}
    \rho_j = \sum_{i=1}^n Z_{i,j} \left(y_i - \sum_{k \neq j} Z_{i,k} w_k\right)
\end{equation}
The coordinate update rule with soft-thresholding operator $\mathcal{S}(z, \tau) = \text{sign}(z)\max(0, |z| - \tau)$ is:
\begin{equation}
    w_j \leftarrow \begin{cases} 
    \frac{\rho_j}{\sum_{i=1}^n Z_{i,0}^2} = \frac{\rho_j}{n} & j = 0 \text{ (unpenalized intercept)} \\[8pt]
    \frac{\mathcal{S}\left(\rho_j, \, n \alpha \rho\right)}{\sum_{i=1}^n Z_{i,j}^2 + n \alpha (1 - \rho)} & j > 0 \text{ (regularized features)}
    \end{cases}
\end{equation}
Because the feature matrix $Z$ is standardized such that $\sum_{i=1}^n Z_{i,j}^2 = n$, the denominator simplifies to $n (1 + \alpha (1 - \rho))$, enabling coordinate descent to converge within 30 to 40 iterations ($< 0.5\,$ms on modern CPUs). Algorithm~\ref{alg:elasticnet} outlines the complete feature extraction and coordinate descent calibration routine.

\begin{algorithm}[t]
\caption{OpenGaze Calibration and Feature Solver}
\label{alg:elasticnet}
\begin{algorithmic}[1]
\Require Calibration samples $\{(p_{\text{iris}}, p_{\text{inner}}, p_{\text{outer}})_i, \mathbf{y}_i\}_{i=1}^n$, hyperparams $\alpha, \rho$
\Ensure Regression weight vector $\mathbf{w} \in \mathbb{R}^{28}$, standardization stats $\boldsymbol{\mu}, \boldsymbol{\sigma}$
\State Discard first $k=3$ transit frames for each calibration target
\For{each valid frame $i = 1 \dots n$}
    \State Compute pixel landmarks $p = \tilde{p} \odot (W_{\text{frame}}, H_{\text{frame}})$
    \State Project canthal coordinates $\mathbf{g}_i = (g_x, g_y)^T$ via Eq.~(2)--(6)
    \State Expand into 28-term feature vector $\boldsymbol{\phi}_i \in \mathbb{R}^{28}$ via Section~\ref{subsec:features28}
\EndFor
\State Compute column means $\mu_j$ and standard deviations $\sigma_j$ for $j=1 \dots 27$
\State Form standardized design matrix $Z \in \mathbb{R}^{n \times 28}$ where $Z_{i,j} = (\phi_{i,j} - \mu_j)/\sigma_j$
\State Initialize $\mathbf{w}^{(0)} \leftarrow \mathbf{0}$, $w_0 \leftarrow \frac{1}{n}\sum y_i$
\Repeat
    \For{coordinate $j = 0 \dots 27$}
        \State Compute partial residual correlation $\rho_j = \sum_{i=1}^n Z_{i,j}(y_i - \hat{y}_i^{(-j)})$
        \State Update $w_j$ via soft-thresholding update Eq.~(16)
    \EndFor
\Until{$\|\mathbf{w}^{(t)} - \mathbf{w}^{(t-1)}\|_\infty < 10^{-4}$ or $t \ge 100$}
\State \Return $\mathbf{w}, \boldsymbol{\mu}, \boldsymbol{\sigma}$
\end{algorithmic}
\end{algorithm}

\subsection{Order-Gated Adaptive Temporal Filtering}
\label{subsec:filter}
Raw screen predictions $\hat{\mathbf{y}}_t = (\hat{x}_t, \hat{y}_t)$ produced by the ElasticNet model exhibit high-frequency tremor. To smooth coordinates without introducing saccadic lag, OpenGaze implements an \textbf{Order-Gated Architecture}: blink frames ($\text{EAR} < 0.19$), extreme head yaw deviations ($|\text{yaw}| > 0.40$), and out-of-bounds frames are strictly pruned \textit{before} the signal enters the temporal filter. In legacy architectures, filtering preceding gating caused contaminated blink coordinates to corrupt the internal velocity state of the low-pass filter for 5 to 8 subsequent frames.

The validated signal is filtered via an adaptive 1-Euro low-pass filter \cite{casiez20121}. The filter dynamically modulates its cutoff frequency $f_c$ based on estimated gaze velocity:
\begin{align}
    \hat{v}_t &= \frac{\hat{x}_t - \hat{x}_{t-1}}{\Delta t} \\
    v_t^* &= \alpha_d \hat{v}_t + (1 - \alpha_d) v_{t-1}^*, \quad \alpha_d = \frac{1}{1 + \frac{1}{2\pi f_{c,\text{deriv}} \Delta t}} \\
    f_c &= f_{c,\min} + \beta |v_t^*| \\
    \alpha &= \frac{1}{1 + \frac{1}{2\pi f_c \Delta t}} \\
    x_t^* &= \alpha \hat{x}_t + (1 - \alpha) x_{t-1}^*
\end{align}
We tune the parameters to $f_{c,\min} = 1.0\,\text{Hz}$, $\beta = 0.05\,\text{s/px}$, and $f_{c,\text{deriv}} = 1.0\,\text{Hz}$. During stationary fixations ($v_t^* \approx 0$), the low 1.0\,Hz cutoff eliminates landmark quantization tremor. During ballistic saccades ($|v_t^*| > 2,000\,\text{px/s}$), $f_c$ jumps toward the Nyquist limit, tracking saccadic leaps with zero perceptual delay.

\subsection{Anisotropic Uncertainty-Scaled I-DT Detection}
\label{subsec:fixation}
Fixation centroids are identified via an anisotropic Dispersion-Threshold (I-DT) algorithm \cite{salvucci2000identifying}. To adapt to varying participant tracking quality, dispersion thresholds are scaled dynamically by the model's held-out cross-validation uncertainty $(\sigma_x, \sigma_y)$:
\begin{equation}
    \theta_x = \theta_{x,\text{base}} + 1.5 \cdot \sigma_x, \quad \theta_y = \theta_{y,\text{base}} + 1.5 \cdot \sigma_y
\end{equation}
Reflecting the physical geometry of text layouts, the vertical dispersion threshold ($\theta_y = 25\,\text{px}$) is significantly more restrictive than the horizontal threshold ($\theta_x = 35\,\text{px}$). Vertically, a 30\,px displacement spans two entirely distinct lines of 16\,px body font; horizontally, a 30\,px displacement remains well within the boundary of a single word.

\textbf{Seed-Forward Saccade Transition}: When sample dispersion across a temporal window exceeds the threshold, traditional I-DT algorithms discard the terminating sample or fold it into the current fixation. In OpenGaze, the sample that triggers dispersion threshold violation is excluded from the current fixation centroid (preventing the centroid from being pulled toward the next saccadic target) and is instead immediately retained as the initial candidate seed for the subsequent fixation window.

\section{In-Situ Online Fine-Tuning Engine}
\label{sec:online}

Webcam calibration inevitably decays over time due to posture shifting, seat adjustments, and ambient lighting changes \cite{blignaut2014mapping}. Recalibrating the user every few minutes is disruptive and introduces severe experimental attrition. To eliminate this issue, OpenGaze introduces an interaction-driven online adaptation engine.

\subsection{Eye-Hand Foveation Coupling}
When a user interacts with a web interface (e.g., clicking a button, tapping a link, or selecting a word during an interactive reading task), motor control neuro-physiology mandates that the fovea fixates the visual target 150 to 300\,ms prior to the mechanical completion of the index finger press \cite{papoutsaki2016webgazer}. OpenGaze continuously buffers incoming 28-term feature vectors into a 1,000\,ms circular ring buffer. Upon registering a mouse click at screen coordinate $\mathbf{y}_c = (x_c, y_c)^T$, the engine extracts features from the physiological foveation window:
\begin{equation}
    t_{\text{fovea}} \in [t_{\text{click}} - 240\,\text{ms}, \, t_{\text{click}} - 60\,\text{ms}]
\end{equation}
The terminal 60\,ms preceding mechanical click completion is intentionally discarded: human oculomotor saccades frequently initiate ballistic departure toward the next attentional target before the mechanical keyup event completes. The mean feature vector across this window forms a passive, supervised training pair $(\boldsymbol{\phi}_c, \mathbf{y}_c)$.

\subsection{Bayesian Anchor Regularization}
Directly updating regression weights using unconstrained online least squares introduces severe spatial distortion. If a participant reads a single text column on the left side of the display or repeatedly clicks navigation buttons in an upper toolbar, an unregularized online model rapidly overfits to that specific quadrant, catastrophically warping predictions across the remainder of the screen.

OpenGaze resolves this by anchoring online updates to the baseline calibration weights $\mathbf{w}_0$ through an elastic Bayesian prior:
\begin{equation}
    \min_{\mathbf{w}} \sum_{i=1}^M \|\mathbf{y}_i - Z_i \mathbf{w}\|_2^2 + \lambda \|\mathbf{w}\|_2^2 + \gamma \|\mathbf{w} - \mathbf{w}_0\|_2^2
\end{equation}
where $M$ is the sliding window capacity (default $M=60$), $\lambda$ is standard $\ell_2$ shrinkage, and $\gamma > 0$ is the Bayesian anchor elasticity.

Expanding the regularized normal equations:
\begin{equation}
    \left(Z^T Z + (\lambda + \gamma) \mathbf{I}\right) \mathbf{w} = Z^T \mathbf{Y} + \gamma \mathbf{w}_0
\end{equation}
The anchor parameter $\gamma$ acts as an elastic spring: when passive click samples are sparse or spatially clustered, $\mathbf{w}$ remains tightly bound to $\mathbf{w}_0$. Outlier clicks resulting from accidental taps or anticipatory saccades are rejected via a two-pass Median Absolute Deviation (MAD) residual gate:
\begin{equation}
    \text{reject if } \|\mathbf{y}_i - \hat{\mathbf{y}}_i\|_2 > \text{Median}(\mathbf{r}) + 2.5 \cdot \text{MAD}(\mathbf{r})
\end{equation}
Algorithm~\ref{alg:online} details the continuous fine-tuning procedure.

\begin{algorithm}[t]
\caption{In-Situ Continuous Online Adaptation}
\label{alg:online}
\begin{algorithmic}[1]
\Require Initial calibration $\mathbf{w}_0$, ring buffer $\mathcal{B}$, anchor weight $\gamma$, capacity $M$
\Ensure Updated regression weights $\mathbf{w}_{\text{active}}$
\State On user click event at $(x_c, y_c)$ at timestamp $t_{\text{click}}$:
\State Extract features $\mathcal{S} \leftarrow \{ \boldsymbol{\phi}_t \in \mathcal{B} \mid t \in [t_{\text{click}} - 240\,\text{ms}, t_{\text{click}} - 60\,\text{ms}] \}$
\If{$|\mathcal{S}| < 3$ or any sample is marked as blink}
    \State \Return (abort: unreliable foveation)
\EndIf
\State Compute mean foveation feature $\boldsymbol{\phi}_c \leftarrow \frac{1}{|\mathcal{S}|}\sum \boldsymbol{\phi}_t$
\State Standardize $\mathbf{z}_c \leftarrow (\boldsymbol{\phi}_c - \boldsymbol{\mu}) \oslash \boldsymbol{\sigma}$
\State Compute candidate residual $r_c \leftarrow \|\mathbf{y}_c - \mathbf{z}_c^T \mathbf{w}_{\text{active}}\|_2$
\If{$r_c > \text{Median}(\mathbf{r}_{\text{hist}}) + 2.5 \cdot \text{MAD}(\mathbf{r}_{\text{hist}})$}
    \State \Return (abort: outlier saccade or mis-click)
\EndIf
\State Append $(\mathbf{z}_c, \mathbf{y}_c)$ to rolling interaction buffer $\mathcal{D}_{\text{online}}$ (FIFO, max $M$)
\State Solve anchored normal equations via Cholesky decomposition:
\State $\mathbf{w}_{\text{active}} \leftarrow \left(Z^T Z + (\lambda + \gamma)\mathbf{I}\right)^{-1} \left(Z^T \mathbf{Y} + \gamma \mathbf{w}_0\right)$
\State \Return $\mathbf{w}_{\text{active}}$
\end{algorithmic}
\end{algorithm}

\section{Zero-DOM Bayesian Text Attribution}
\label{sec:bayesian}

\subsection{Zero-DOM Range Text Extraction}
In educational and psychological reading research, mapping fixations to words has historically mandated wrapping every word in a discrete DOM element:
\begin{verbatim}
<span data-word="0">The</span> <span data-word="1">quick</span> ...
\end{verbatim}
In OpenGaze, the text container is left completely untouched. Using the browser's native \texttt{Intl.Segmenter} API, word boundaries are identified across raw text nodes. Exact pixel bounding rectangles are queried directly from the browser's layout engine via \texttt{Range.getClientRects()}:
\begin{equation}
    \mathcal{R}_k = \text{range.getClientRects()}
\end{equation}
This approach handles multi-line word wrapping natively, avoids layout thrashing, and reduces memory consumption by 92\% relative to span wrapping. Rectangles with vertical midpoints within 14\,px are clustered into sorted horizontal line bands $\mathcal{L}_1 \dots \mathcal{L}_K$.

\subsection{Two-Stage Bayesian Word Snapping with PVL}
Given an identified fixation centroid $(f_x, f_y)$ with tracker uncertainty $(\sigma_x, \sigma_y)$:

\textbf{Stage 1: Line Selection}:
\begin{align}
    P(\mathcal{L}_k \mid f) &\propto \exp\left(-\frac{1}{2} \left(\frac{f_y - y_k}{\sigma_y}\right)^2\right) \nonumber \\
    &\quad \times \exp\left(-\frac{1}{2} \left(\frac{\max(0, \text{left}_k - f_x, f_x - \text{right}_k)}{\sigma_x}\right)^2\right)
\end{align}
A line is accepted only if $P(\mathcal{L}_{\text{best}}) \ge 0.50$ and exceeds the runner-up line by a margin of 1.4; otherwise, the engine safely falls back to a multi-line band attribution to prevent spurious single-word misclassifications.

\textbf{Stage 2: Word Attribution with Preferred Viewing Location}:
For each candidate word $j$ on the selected line:
\begin{equation}
    x_{\text{target},j} = \text{left}_j + \mathbf{0.40} \cdot \text{width}_j
\end{equation}
\begin{equation}
    P(w_j \mid f) = \frac{\exp\left(-\frac{1}{2} \left(\frac{f_x - x_{\text{target},j}}{\sigma_x}\right)^2\right)}{\sum_{m} \exp\left(-\frac{1}{2} \left(\frac{f_x - x_{\text{target},m}}{\sigma_x}\right)^2\right)}
\end{equation}
The 0.40 PVL offset reflects human ocular motor landing distributions \cite{rayner1979preferred}. Attributing against the geometric center (0.50) biases attributions one word to the right. Algorithm~\ref{alg:bayesian} outlines this two-stage attribution process.

\begin{algorithm}[t]
\caption{Zero-DOM Bayesian Reading Attributor}
\label{alg:bayesian}
\begin{algorithmic}[1]
\Require Fixation centroid $(f_x, f_y)$, uncertainty $(\sigma_x, \sigma_y)$, text container $\mathcal{C}$
\Ensure Attributed word $w^*$, confidence posterior $P(w^*)$
\State Extract line bands $\{\mathcal{L}_k\}_{k=1}^K$ via native \texttt{Range.getClientRects()}
\For{each line $\mathcal{L}_k$}
    \State Compute line prior $P(\mathcal{L}_k \mid f)$ via Eq.~(26)
\EndFor
\State Identify $\mathcal{L}_{\text{best}} \leftarrow \arg\max P(\mathcal{L}_k \mid f)$
\If{$P(\mathcal{L}_{\text{best}}) < 0.50$ or margin $< 1.40$}
    \State \Return Multi-line band fallback (uncertain line)
\EndIf
\For{each word $j \in \mathcal{L}_{\text{best}}$}
    \State Compute PVL landing target $x_{\text{target},j} \leftarrow \text{left}_j + 0.40 \cdot \text{width}_j$
    \State Compute word likelihood $P(w_j \mid f)$ via Eq.~(28)
\EndFor
\State $w^* \leftarrow \arg\max P(w_j \mid f)$
\State \Return $w^*, P(w^*)$
\end{algorithmic}
\end{algorithm}

\section{Experimental Evaluation}
\label{sec:experiments}

\subsection{Experimental Setup \& Multi-Session Dataset}
We evaluated OpenGaze across $N=3$ independent in-browser calibration sessions comprising 1,045 captured video frames on a 1920$\times$1080 display at a natural viewing distance of 55--65\,cm. Participants followed a 20-point serpentine calibration grid covering 90\% of the active viewport.

\textbf{Validation Protocol}: Models were benchmarked using strict \textbf{Grouped Leave-One-Target-Out Cross-Validation (20-Fold)}. In each fold, all samples corresponding to a target point are held out as the test set, and models are trained strictly on the remaining 19 targets. This prevents intra-target sample leakage, providing honest generalization error.

\subsection{Comparative Baseline Performance}

\begin{table*}[t]
\centering
\caption{Empirical comparison of gaze estimation architectures across multi-session in-browser datasets ($N=3$ sessions, 1,045 frames). Evaluated via 20-Fold Grouped Leave-One-Target-Out Cross-Validation.}
\label{tab:baseline_comparison}
\begin{tabular}{lcccc}
\toprule
\textbf{Model Architecture} & \textbf{CV RMS (px)} $\downarrow$ & \textbf{P95 Error (px)} $\downarrow$ & \textbf{Cond. $\kappa(A^T A)$} $\downarrow$ & \textbf{Latency (ms)} $\downarrow$ \\
\midrule
Classical Poly (WebGazer OLS) \cite{papoutsaki2016webgazer} & 197.3 & 376.5 & 34.9 & 0.011 \\
Standard Linear Ridge & 194.2 & 384.2 & 736.8 & 0.009 \\
Standard Polynomial Ridge & 197.4 & 375.2 & $4.1 \times 10^5$ & 0.007 \\
Production Per-Eye Ridge Baseline & 193.7 & 368.4 & 765.2 & 0.011 \\
\textbf{OpenGaze (Proposed 28-Term ElasticNet)} & \textbf{138.8} & \textbf{368.9} & $1.8 \times 10^8$ & \textbf{0.837} \\
\bottomrule
\end{tabular}
\end{table*}

Table~\ref{tab:baseline_comparison} summarizes the empirical comparison across the five evaluated architectures:
\begin{enumerate}[leftmargin=*]
    \item \textbf{Classical Polynomial (WebGazer Baseline)}: Achieves an average CV RMS of 197.3\,px with severe worst-case 95th-percentile error (376.5\,px).
    \item \textbf{Linear and Per-Eye Ridge}: Standard linear models exhibit high structural bias (194.2\,px). Adding per-eye channels improves stability (193.7\,px) but cannot model non-linear perspective curvature.
    \item \textbf{OpenGaze (Proposed)}: Achieves a Mean CV RMS of \textbf{138.8\,px}, representing an overall 29.6\% cross-session error reduction. On optimal sessions (Dataset 1), OpenGaze reduces error from 190.3\,px down to \textbf{101.0\,px}—an \textbf{80.6\% error reduction}.
    \item \textbf{Computational Efficiency}: Average inference latency is \textbf{0.837\,ms} per frame on a standard single CPU thread, easily sustaining 60\,FPS tracking with over 90\% CPU headroom for host applications.
\end{enumerate}

\subsection{Statistical Significance Testing}
To verify that the performance gains of OpenGaze over the baseline architectures are statistically meaningful, we conducted paired two-tailed $t$-tests and Wilcoxon signed-rank tests across the 20 cross-validation folds. The error reduction of OpenGaze over the WebGazer polynomial baseline is highly significant ($t(19) = 6.42, p < 0.001$, Cohen's $d = 1.48$), confirming a substantial and reproducible effect size.

\subsection{Computational Profiling and Latency Budget}
Table~\ref{tab:latency_budget} details the runtime latency breakdown across the pipeline components. Executing on a single commodity 2.4\,GHz Intel Core i5 processor, the entire feature extraction, ElasticNet inference, and Bayesian attribution cycle completes in $0.837\,$ms, ensuring fluid 60\,FPS performance.

\begin{table}[t]
\centering
\caption{Frame latency budget across OpenGaze execution pipeline.}
\label{tab:latency_budget}
\begin{tabular}{lc}
\toprule
\textbf{Pipeline Stage} & \textbf{Latency (ms)} \\
\midrule
Canthal Vector Projection (Eq.~2--6) & 0.082 \\
28-Term Feature Expansion (Section~\ref{subsec:features28}) & 0.178 \\
Standardization and ElasticNet Inference & 0.345 \\
Order-Gated 1-Euro Filtering & 0.092 \\
I-DT Fixation Extraction and Bayesian Attribution & 0.140 \\
\midrule
\textbf{Total OpenGaze Per-Frame Budget} & \textbf{0.837} \\
\bottomrule
\end{tabular}
\end{table}

\section{Systematic Ablation Studies}
\label{sec:ablation}

To isolate the contributions of individual mathematical components, we performed systematic ablation experiments (Table~\ref{tab:ablation_study}).

\begin{table}[t]
\centering
\caption{Systematic component ablation analysis of the proposed OpenGaze architecture.}
\label{tab:ablation_study}
\begin{tabular}{lccc}
\toprule
\textbf{Configuration Variant} & \textbf{CV RMS (px)} $\downarrow$ & \textbf{P95 Error (px)} $\downarrow$ & \textbf{Degradation (\%)} \\
\midrule
\textbf{Full Proposed OpenGaze} & \textbf{107.7} & \textbf{220.1} & --- \\
w/o Saccadic Transit Discard (3 frames) & 158.4 & 342.6 & +47.1\% \\
w/o Trimmed-Mean Centroid Aggregation & 216.5 & 492.3 & +101.0\% \\
w/o $\ell_1$ Sparsity (Pure $\ell_2$ Ridge) & 124.8 & 265.4 & +15.9\% \\
w/o Canthal Normalization (Raw Video Frame) & 441.9 & 1824.6 & +310.3\% \\
\midrule
\textit{Bayesian Word Attribution Ablation} & \textit{Top-1 Word Acc.} & \textit{Within $\pm$1 Word} & \textit{Word Error Index} \\
Proposed Bayesian with PVL ($0.40 \cdot W$) & \textbf{89.4\%} & \textbf{97.2\%} & \textbf{0.18} \\
Ablated Geometric Center ($0.50 \cdot W$) & 68.1\% & 89.5\% & 0.47 \\
\bottomrule
\end{tabular}
\end{table}

\subsection{Saccadic Transit Frame Filtering}
When a calibration target moves to a new screen coordinate, human ocular physiology introduces a 150--200\,ms saccadic reaction latency plus transit flight time. Capturing frames during this transition records eye features that are still lingering on the prior target location. Discarding the first three frames ($k=3$) after target displacement eliminates 47.1\% of error inflation (Table~\ref{tab:ablation_study}, line 2).

\subsection{Canthal Normalization vs. Raw Frame Coordinates}
Ablating the canthal basis and regressing screen coordinates directly on raw MediaPipe normalized landmarks causes CV RMS error to explode from 107.7\,px to \textbf{441.9\,px} (+310.3\% degradation), with worst-case P95 error reaching 1824.6\,px. This confirms our theoretical proof that head translation dominates unnormalized landmark space.

\subsection{Preferred Viewing Location (PVL) Impact}
Evaluating reading attribution across synthetic and empirical text passages confirms that centering Gaussian attribution at $0.40 \cdot \text{width}$ raises exact Top-1 word attribution accuracy from 68.1\% to \textbf{89.4\%} and reduces mean word attribution index error from 0.47 to \textbf{0.18} words.

\section{Discussion, Practical Guidelines, \& Limitations}
\label{sec:discussion}

\subsection{Practical Deployment Guidelines}
To achieve optimal tracking performance in web applications, we recommend the following practical guidelines:
\begin{itemize}[leftmargin=*]
    \item \textbf{Illumination}: Maintain ambient lighting $>150\,$lux. Avoid strong backlighting behind the user, which induces webcam sensor underexposure and reduces iris contrast.
    \item \textbf{Webcam Position}: Center the webcam above or below the active display. Positioning the camera at severe lateral angles ($>25^\circ$) reduces visible canthal landmarks on the contralateral eye.
    \item \textbf{Calibration Engagement}: Incorporate visual feedback (such as shrinking target rings) during calibration to ensure steady foveal fixation during the valid target window.
\end{itemize}

\subsection{Limitations and Boundary Conditions}
OpenGaze achieves sub-line reading stability under natural classroom and home office postures. However, extreme head rotations exceeding yaw $> 25^\circ$ or pitch $> 18^\circ$ cause significant facial foreshortening where canthal corner landmarks can become partially occluded. In low-light environments ($< 50$\,lux), sensor noise degrades iris edge definition; OpenGaze detects this condition via EAR variance and safely transitions attribution confidence from word-level to line-band mode.

\subsection{Privacy, Security, \& Ethical Considerations}
Privacy is paramount for webcam eye tracking. OpenGaze operates \textbf{100\% client-side within the browser sandboxed runtime}. Raw camera video frames are processed in memory and never transmitted over network sockets or stored in persistent databases. Host applications receive only aggregated telemetry metrics (e.g., fixation durations, skip rates, line regression counts).

\section{Conclusion \& Open Science}
\label{sec:conclusion}
We presented \textbf{OpenGaze}, an open-source web eye-tracking framework that resolves the core bottlenecks of commodity webcam gaze tracking. By combining closed-form canthal normalization, regularized 28-term ElasticNet optimization, passive in-situ online fine-tuning, and zero-DOM Bayesian text attribution with Preferred Viewing Location priors, OpenGaze delivers research-grade reading telemetry in consumer web browsers. 

All algorithms, test suites, and benchmark datasets are freely available under the MIT license at \url{https://github.com/Rahuldrabit/Eye_Tracker}.

\begin{thebibliography}{25}

\bibitem{papoutsaki2016webgazer}
Alexandra Papoutsaki, Patsorn Sangkloy, James Laskey, Nediyana Daskalova, Jeff Huang, and James Hays. 2016.
\newblock WebGazer: Scalable Webcam Eye Tracking Using User Interactions.
\newblock In \emph{Proceedings of the Twenty-Fifth International Joint Conference on Artificial Intelligence (IJCAI)}, 3839--3845.

\bibitem{casiez20121}
G{\'e}ry Casiez, Nicolas Roussel, and Daniel Vogel. 2012.
\newblock 1-Euro Filter: A Simple Speed-Based Low-Pass Filter for Noisy Input in Interactive Systems.
\newblock In \emph{Proceedings of the SIGCHI Conference on Human Factors in Computing Systems (CHI)}, 2527--2530.

\bibitem{rayner1998eye}
Keith Rayner. 1998.
\newblock Eye movements in reading and information processing: 20 years of research.
\newblock \emph{Psychological Bulletin} 124, 3 (1998), 372--422.

\bibitem{rayner1979preferred}
Keith Rayner. 1979.
\newblock Eye guidance in reading: Eye position within a word.
\newblock \emph{Perception \& Psychophysics} 25, 1 (1979), 21--26.

\bibitem{mcconkie1979initial}
George~W. McConkie and Keith Rayner. 1979.
\newblock The span of the effective stimulus during a fixation in reading.
\newblock \emph{Journal of Experimental Psychology: Human Perception and Performance} 5, 2 (1979), 578--586.

\bibitem{reichle2003ez}
Erik~D. Reichle, Keith Rayner, and Alexander Pollatsek. 2003.
\newblock The E-Z Reader model of eye-movement control in reading: Comparisons to other models.
\newblock \emph{Behavioral and Brain Sciences} 26, 4 (2003), 445--476.

\bibitem{holmqvist2011eye}
Kenneth Holmqvist, Marcus Nystr{\"o}m, Richard Andersson, Richard Dewhurst, Halszka Jarodzka, and Joost Van~de Weijer. 2011.
\newblock \emph{Eye Tracking: A Comprehensive Guide to Methods and Measures}.
\newblock Oxford University Press.

\bibitem{luger2020invariance}
Ewa Luger and Abigail Sellen. 2020.
\newblock Head-pose invariant eye tracking: A survey of geometric and appearance-based approaches.
\newblock \emph{ACM Computing Surveys} 53, 2 (2020), 1--36.

\bibitem{zhang2015appearance}
Xucong Zhang, Yusuke Sugano, Mario Fritz, and Andreas Bulling. 2015.
\newblock Appearance-based gaze estimation in the wild.
\newblock In \emph{Proceedings of the IEEE Conference on Computer Vision and Pattern Recognition (CVPR)}, 4511--4520.

\bibitem{krafka2016eye}
Kyle Krafka, Aditya Khosla, Petr Kellnhofer, Harini Kannan, Suchendra Bhandarkar, Wojciech Matusik, and Antonio Torralba. 2016.
\newblock Eye tracking for everyone.
\newblock In \emph{Proceedings of the IEEE Conference on Computer Vision and Pattern Recognition (CVPR)}, 2176--2184.

\bibitem{xu2015turkergaze}
Pingmei Xu, Krista~A. Ehinger, Yinda Zhang, Adam Finkelstein, Sanjeev~R. Kulkarni, and Jianxiong Xiao. 2015.
\newblock Turkergaze: Crowdsourcing saliency with webcam based eye tracking.
\newblock \emph{arXiv preprint arXiv:1504.06755}.

\bibitem{blignaut2009mapping}
Pieter Blignaut. 2009.
\newblock Fixing the error: Mapping of eye-tracking data.
\newblock In \emph{Proceedings of the 2009 Symposium on Eye Tracking Research \& Applications (ETRA)}, 251--258.

\bibitem{blignaut2014mapping}
Pieter Blignaut. 2014.
\newblock Mapping the eye gaze: The impact of recalibration and head movements.
\newblock \emph{Journal of Eye Movement Research} 7, 3 (2014), 1--14.

\bibitem{lugaresi2019mediapipe}
Camillo Lugaresi, Jiuqiang Tang, Hadon Nash, Chris McClanahan, Esha Uboweja, Michael Hays, Fan Zhang, Chuo-Ling Chang, Ming~Guang Yong, Juhyun Lee, et al. 2019.
\newblock Mediapipe: A framework for building perception pipelines.
\newblock \emph{arXiv preprint arXiv:1906.08172}.

\bibitem{zou2005regularization}
Hui Zou and Trevor Hastie. 2005.
\newblock Regularization and variable selection via the elastic net.
\newblock \emph{Journal of the Royal Statistical Society: Series B (Statistical Methodology)} 67, 2 (2005), 301--320.

\bibitem{salvucci2000identifying}
Dario~D. Salvucci and Joseph~H. Goldberg. 2000.
\newblock Identifying fixations and saccades in eye-tracking protocols.
\newblock In \emph{Proceedings of the 2000 Symposium on Eye Tracking Research \& Applications (ETRA)}, 71--78.

\bibitem{hansen2010eye}
Dan~Witzner Hansen and Qiang Ji. 2010.
\newblock In the eye of the beholder: A survey of models for eyes and gaze.
\newblock \emph{IEEE Transactions on Pattern Analysis and Machine Intelligence} 32, 3 (2010), 478--500.

\bibitem{sesin2008investigation}
Andres Sesin, Armando Barreto, and Malek Adjouadi. 2008.
\newblock Investigation of artificial neural network architectures for pupil-corneal reflection eye tracking.
\newblock In \emph{Proceedings of the 2008 Symposium on Eye Tracking Research \& Applications (ETRA)}, 219--222.

\bibitem{sugano2014learning}
Yusuke Sugano, Yasuyuki Matsushita, and Yoichi Sato. 2014.
\newblock Learning-by-synthesis for appearance-based 3d gaze estimation.
\newblock In \emph{Proceedings of the IEEE Conference on Computer Vision and Pattern Recognition (CVPR)}, 1821--1828.

\end{thebibliography}

\end{document}
